\section{Method}
\label{sec:method}

\subsection{Policies with a known order}
\label{sec:method:popA}
Let $\pi$ be a pretrained policy that maps the sensor input of a scenario to a planned trajectory
of eight waypoints over four seconds.
A degradation family $d$ with levels $\ell\in\{0,1,2,3\}$ defines policies $\pi_{d,\ell}$, where
$\pi_{d,0}=\pi$ and the corruption grows with $\ell$.
For every family the stochastic part (noise field, drift direction) is drawn from a generator keyed
by scenario and family but \emph{not} by level, so the levels of a family differ only in magnitude
(common random numbers); paired differences between levels then reflect the level, not the draw.

We use six families (\cref{tab:degradations}).
Three act on the input: \emph{camera dropout} blanks the cameras the policy reads, front camera first;
\emph{image noise} adds Gaussian pixel noise; \emph{stale observation} serves the frame of
$0.5$--$1.5$\,s ago as the current one.
Three act on the output: \emph{lateral noise} adds a lateral drift that grows linearly along the
plan; \emph{late braking} stretches the plan along its own path; \emph{latency} executes the plan
that was computed $k$ replanning steps earlier.
In closed loop, latency keeps the last $k$ plans in world coordinates and re-expresses the oldest in
the current ego frame, dropping its elapsed waypoints.
In single-shot evaluators, where no earlier plan exists, latency is reproduced exactly as what was
executed: the policy plans from the full input of $k$ frames earlier (sensors and ego state), and the
plan is moved into the current ego frame through the logged ego pose, its first $k$ waypoints dropped
and $k$ appended at the last waypoint's velocity.
The four-frame history of NAVSIM caps $k$ at three frames ($1.5$\,s); the closed-loop level~3 is
$2.0$\,s.

\subsection{Validating the order}
\label{sec:method:mono}
``Higher level is worse'' is an assumption about the world, so we test it with a reference evaluator
$E^\star$ before using it.
For base policy $b$, family $d$ and adjacent levels $\ell,\ell+1$, let
$\Delta_s = E^\star(\pi_{d,\ell},s) - E^\star(\pi_{d,\ell+1},s)$ be the paired score drop on scenario
$s$.
We resample scenarios $B=1000$ times and compute $P_{\ell} = \Pr(\bar\Delta > 0)$.
Family $d$ is \emph{retained} if $P_\ell > 0.95$ for all three steps and for every base policy
(an intersection--union test); otherwise it is discarded and reported, not re-levelled.

\paragraph{Selection on held-out scenarios.}
If degradations are selected on the same scenarios on which evaluators are later compared, the
reference recovers the order of the retained families by construction.
We therefore assign each scenario to a selection or an evaluation half by a hash of its token
(SHA-256, fixed salt, fraction $0.5$), so that a scenario shared by several evaluators falls in the
same half for all of them.
Monotonicity is tested on the selection half; every other quantity, for every evaluator including
$E^\star$, is computed on the evaluation half.

\subsection{Measuring an evaluator}
\label{sec:method:metrics}
For an evaluator $E$ and a retained family, the true order of the four levels is known.
\emph{Ordering recovery} is Kendall's $\tau$ between the true order and the order of $E$'s mean
scores, averaged over base policies, with a bootstrap interval over scenarios, together with the
probability that the full order is recovered.
Because $\tau$ saturates when many scenarios are averaged, we also report \emph{sample efficiency}:
the smallest number of scenarios, on a grid from 10 to 800, at which the full order is recovered in
at least 90\% of subsamples.

\subsection{Specificity: changes that should not count}
\label{sec:method:spec}
An evaluator can recover every degradation order merely by penalising any change to the plan.
Open-loop displacement does exactly that.
To separate sensitivity to bad driving from sensitivity to change, we construct benign twins of the
two output families.
A replanning controller executes only the beginning of each plan.
In our closed-loop protocol the plan is replaced every $0.5$\,s; the speed target is derived from the
first second of waypoints after the replan offset (up to $1.4$\,s of the plan), and steering pursues
the waypoint at distance $4+0.1v$\,m.
We therefore define the \emph{tail} of a plan as everything beyond the first waypoint that is at least
$1.5$\,s and $10$\,m ahead, and apply lateral noise (\emph{tail lateral}) or the path stretch
(\emph{tail stretch}) to the tail only, with the same parameters and drift direction as the harmful
family.
Under this protocol the executed controls are unchanged; we verify this against the controller's
rules for speeds of 1--20\,m/s and curvatures up to $0.05$\,m$^{-1}$, and empirically below.

A benign family is \emph{validated} if, for every base and level, the 90\% bootstrap interval of the
reference's mean drop lies within $\pm0.01$ (one point of the reference score; a TOST equivalence test
at $\alpha=0.05$), on the selection half.
An evaluator \emph{penalises} a validated benign family if, at some base and level, its drop is
positive with probability above $0.95$ and at least 25\% of its drop for the harmful twin at the same
level.
All thresholds were fixed before the benign runs.

\subsection{Policies with no known order}
\label{sec:method:popB}
Degradations test detection; deployment needs ranking of genuinely different policies.
For a natural population (different architectures and training seeds) no order is known, so we
measure agreement: Kendall's $\tau$ between evaluators' policy rankings with a paired bootstrap over
scenarios, family-level contrasts with intervals, and the per-scenario correlation of those
contrasts, which we compare with the closed-loop test-retest correlation obtained by re-running the
same policies.
